\documentclass[11pt,a4paper]{article}

\usepackage[head=2cm, bottom=2.5cm, left=2.5cm, right=2.5cm]{geometry}
\usepackage{amsmath}
\usepackage{amsfonts}
\usepackage{amssymb}
\usepackage{amsthm}
\usepackage{bbm}
\usepackage{booktabs}
\usepackage{array}
\usepackage{enumitem}
\usepackage[hidelinks]{hyperref}

\newtheorem{theorem}{Theorem}
\newtheorem{proposition}{Proposition}
\newtheorem{definition}{Definition}
\newtheorem{example}{Example}
\newtheorem{remark}{Remark}
\newtheorem{lemma}{Lemma}

\newcommand{\VaR}{\operatorname{VaR}}
\newcommand{\ES}{\operatorname{ES}}
\newcommand{\Par}{\operatorname{Par}}
\newcommand{\Exp}{\operatorname{Exp}}
\newcommand{\real}{\mathbb{R}}
\newcommand{\E}{\mathbb{E}}
\renewcommand{\P}{\mathbb{P}}
\newcommand{\N}{\mathcal{N}}
\newcommand{\dd}{\,\mathrm{d}}
\newcommand{\ind}[1]{\mathbbm{1}_{\{#1\}}}

% Marker used for everything that is *not* on the slides: details that were
% left to the (black/white) board, or steps of a proof that were only sketched.
\newcommand{\der}[1]{\medskip\noindent\textbf{Derivation (#1).}\ }
\newcommand{\selfcheck}[1]{\medskip\noindent\textbf{Remark (self-check).}\ \textit{#1}}

\title{\bfseries Statistical Methods for Risk Modelling (MFM5130)\\[2mm]
\large Lecture 1 --- Basics of Quantitative Risk Management\thanks{Summary of the slides of Yang Liu (CUHK-Shenzhen, Fall 2026), based on
A.\,J.\,McNeil, R.\,Frey, P.\,Embrechts (2015), \emph{Quantitative Risk Management}, Princeton University Press,
and on \url{https://qrmtutorial.org/}. All steps that the slides left to the board are written out below and
flagged with \textbf{Derivation}.}}
\author{}
\date{}

\begin{document}
\maketitle

\tableofcontents

\section{Quantitative Risk Management and Market Risk}

\subsection{What is risk?}
Risk is any event or action that may adversely affect an organization's ability to achieve its
objectives and to execute its strategies; it is strongly related to uncertainty/randomness. The
course focuses on \emph{statistical} methods for
\begin{itemize}[nosep]
  \item \textbf{market risk}: the risk of losses in positions arising from movements in market prices;
  \item \textbf{credit risk}: the risk of a counterparty failing to meet its obligations in accordance
        with agreed terms (default);
  \item \textbf{operational risk}: the risk of a loss resulting from inadequate or failed internal
        processes, people, systems, or external events.
\end{itemize}

\subsection{The ABC of QRM and market risk}
Risks are random variables, usually aggregated in a portfolio. We work on a probability space
$(\Omega,\mathcal{F},\P)$; $X$ is a random variable (rv) and $x=X(\omega)$ a realisation
($\omega$ = state of nature).

For market risk we consider
\[
V_t=\text{value of a portfolio }P\text{ at time }t,
\]
where a portfolio $P$ is a collection of financial products (stocks, bonds, derivatives, loans);
determining its value (valuation of financial products) is the hard part. Time $t$ typically means
\emph{today}: $V_t$ is assumed known/observable and is then written $v_t$. We are interested in the
future value $V_{t+1}$ of $P$ (a rv), where ``$t+1$'' means one period ahead, e.g.\ 10 days (market
risk), 1 year (credit and operational risk), 20 years (pension funds).

The \textbf{profit-and-loss (P\&L) variable} of $P$ is
\begin{equation}\label{eq:PL}
  PL_{t+1}=L_{t+1}=-(V_{t+1}-v_t),
\end{equation}
where, by convention, losses (which QRM mainly cares about) live in the right tail of the
distribution function (df) of $L_{t+1}$. Ultimately we want the \emph{loss distribution} of
$L_{t+1}$, or some statistics of it (such as high quantiles). Note that we look at \emph{changes} in
$V_t$ rather than at absolute values --- otherwise we would compare apples and pears.

\paragraph{Risk mapping.} The mapping of risks is a representation
\begin{equation}\label{eq:mapping}
  V_t=f(t,\mathbf{Z}_t),
\end{equation}
where $f$ is a measurable function, the \emph{structure function} of $P$, and $\mathbf{Z}_t$ is a
$d$-dimensional random vector of \emph{risk factors} ($d$ often 1000--10000), e.g.\ (logarithmic)
asset prices or exchange rates. The one-period \emph{risk-factor changes} (or (log-)returns) are
\[
  \mathbf{X}_t=\mathbf{Z}_t-\mathbf{Z}_{t-1},
\]
and finding models for $(\mathbf{X}_t)_t$ is the business of econometrics.

\begin{example}[Two stock price series]
For two historical stock price series $(S_{t,j})_{t}$, $j\in\{1,2\}$ (BMW, Siemens; 1972 days), one
plots $(S_{t,j})_t$ and the log-returns $(X_{t,j}=\log(S_{t,j}/S_{t-1,j}))_t$. The price levels
exhibit trends and non-stationarity, the log-returns look (approximately) stationary with
volatility clustering --- which is why we model the $X$'s and not the $S$'s.
\end{example}

\subsection{From risk factors to the loss}
Plugging the risk-factor changes into \eqref{eq:mapping} gives
\begin{equation}\label{eq:losscalc}
  \begin{aligned}
    L_{t+1}&=-(V_{t+1}-v_t)=-\big(f(t+1,\mathbf{Z}_{t+1})-f(t,\mathbf{z}_t)\big)\\
            &=-\big(f(t+1,\mathbf{z}_t+\mathbf{X}_{t+1})-f(t,\mathbf{z}_t)\big)
             =\ell_{[t]}(\mathbf{X}_{t+1}),
  \end{aligned}
\end{equation}
where $\ell_{[t]}(\cdot)$ is the \emph{loss operator} (which banks have to compute/estimate on a
daily basis). The loss distribution is thus determined by the distribution of $\mathbf{X}_{t+1}$;
note that $\ell_{[t]}(\cdot)$ can be highly nonlinear.

\begin{example}[Portfolio of $d$ stocks]\label{ex:stocks}
Let $P$ consist of $d$ stocks $S_{t,1},\dots,S_{t,d}$, where $S_{t,j}$ is the value of stock $j$ at
time $t$ and $\beta_j$ is the number of shares of stock $j$. The structure function is linear,
\[
  V_t=f(t,\mathbf{Z}_t)=\sum_{j=1}^d \beta_j S_{t,j}
      =\sum_{j=1}^d \beta_j \exp(Z_{t,j}),
  \qquad Z_{t,j}=\log S_{t,j},
\]
and the risk-factor changes are the log-returns
$X_{t+1,j}=\log(S_{t+1,j}/S_{t,j})$. Hence
\begin{equation}\label{eq:stockPL}
  \begin{aligned}
    L_{t+1}&=-(V_{t+1}-v_t)
      =\sum_{j=1}^d \beta_j\big(\exp(Z_{t+1,j})-\exp(z_{t,j})\big)\\
      &=\sum_{j=1}^d \beta_j\big(\exp(z_{t,j}+X_{t+1,j})-\exp(z_{t,j})\big)
      =\sum_{j=1}^d \tilde w_{t,j}\big(\exp(X_{t+1,j})-1\big),
  \end{aligned}
\end{equation}
with $\tilde w_{t,j}=\beta_j s_{t,j}$ and $w_{t,j}=\tilde w_{t,j}/v_t$ the relative value of stock
$j$ in $P$. The loss operator is therefore
\begin{equation}\label{eq:stocklossop}
  \ell_{[t]}(\mathbf{x})=\sum_{j=1}^d \tilde w_{t,j}\big(\exp(x_j)-1\big).
\end{equation}
\end{example}

\paragraph{Key task of QRM.} Based on models and statistical estimates (data) for
$(\mathbf{X}_{t+1})_t$, find the df (or a risk measure) of $L_{t+1}$. There are several
possibilities for ``the df''.

\subsection{Conditional versus unconditional loss distribution}
\begin{enumerate}[label=\arabic*)]
  \item \textbf{Conditional} distribution of $L_{t+1}$ given the history up to and including time $t$:
        \[
          F_{L_{t+1}\mid\mathcal{F}_t}(x)=\P(L_{t+1}\le x\mid\mathcal{F}_t),
        \]
        where $(\mathcal{F}_t)_t$ is the filtration of the underlying process
        $(\mathbf{X}_t)_t$: roughly, $\mathcal{F}_t$ contains the historical information of the
        risk-factor changes up to and including time $t$. We then ask for the distribution of
        $L_{t+1}=\ell_{[t]}(\mathbf{X}_{t+1})$ given $\mathcal{F}_t$, i.e.\ the conditional
        (predictive) loss distribution for the next interval. This forces us to model the dynamics
        of the risk-factor time series and, in particular, to predict volatility --- the most
        suitable approach to market risk.
  \item \textbf{Unconditional} distribution, assuming $(\mathbf{X}_t)_t$ is a stationary time
        series:
        \[
          F_{L_{t+1}}(x)=\P(L_{t+1}\le x)
        \]
        for a generic risk-factor-change vector $\mathbf{X}$ with the same distribution
        $F_{\mathbf{X}}$ as $\mathbf{X}_t,\mathbf{X}_{t-1},\dots$. Neglecting the dynamics
        inevitably leads to this view; it is typical in credit risk and sensible for long horizons.
\end{enumerate}
The issue only arises because of the time-series properties of $(\mathbf{X}_t)_{t\in\mathbb N}$: if
the $\mathbf{X}_t$ are iid random vectors it does not arise, but for a strongly stationary
multivariate time series the two cases must be distinguished. If the risk-factor changes cannot be
assumed stationary for at least some window of time reaching from the present into the intermediate
past, any statistical analysis of the loss distribution is difficult. We shall mainly work with the
\textbf{unconditional} case --- typically the simpler approach (``reality'' lies somewhere in
between).

\subsection{Linearization of the loss operator}
Recall \eqref{eq:losscalc}: $L_{t+1}= -\big(f(t+1,\mathbf{z}_t+\mathbf{X}_{t+1})-f(t,\mathbf{z}_t)\big)
=\ell_{[t]}(\mathbf{X}_{t+1})$. Assume the structure function $f$ is differentiable and use a
first-order (Taylor) approximation: the \emph{linearized loss} $L^{\Delta}_{t+1}$ is
\begin{equation}\label{eq:linloss}
  L_{t+1}\;\leftarrow\;L^{\Delta}_{t+1}
   =-\Big(f_t(t,\mathbf{z}_t)+\sum_{j=1}^d f_{z_j}(t,\mathbf{z}_t)\,X_{t+1,j}\Big),
\end{equation}
where $f_t=\partial f/\partial t$ (the mapping explicitly features time only for derivative
portfolios; the term $f_t(t,\mathbf{z}_t)$ is then sometimes neglected) and
$f_{z_j}=\partial f/\partial z_j$, $j\in\{1,\dots,d\}$, are the partial derivatives with respect to
the $j$-th risk factor. The \emph{linearized loss operator} is accordingly
\begin{equation}\label{eq:linlossop}
  \ell_{[t]}(\mathbf{x})\;\approx\;\ell^{\Delta}_{[t]}(\mathbf{x})
   =-\Big(f_t(t,\mathbf{z}_t)+\sum_{j=1}^d f_{z_j}(t,\mathbf{z}_t)\,x_j\Big).
\end{equation}
Linearization is convenient since linear functions of the risk-factor changes are much easier to
handle analytically --- crucial for the \emph{variance--covariance method}. The approximation is
good over short horizons and for portfolios that are nearly linear in the risk-factor changes, and
typically poor under extremes (e.g.\ extreme market behaviour).

\begin{example}[Example \ref{ex:stocks} revisited]
There is no explicit time dependence in the mapping $V_t=\sum_j\beta_j\exp(z_{t,j})$, and
\[
  f_{z_j}(t,\mathbf{z}_t)=\beta_j\exp(z_{t,j}),\qquad j\in\{1,\dots,d\}.
\]
Hence the linearized loss \eqref{eq:linloss} is
\[
  L^{\Delta}_{t+1}=\sum_{j=1}^d \beta_j \exp(z_{t,j})\,X_{t+1,j}
                  =\sum_{j=1}^d \tilde w_{t,j} X_{t+1,j},
\]
which should be compared with the exact formula \eqref{eq:stockPL}.
\der{the linearized loss for a call option}
We only need to check the derivatives, since \eqref{eq:linloss} is just Taylor's formula:
\[
  \frac{\partial}{\partial x_j}\Big[\tilde w_{t,j}\big(e^{x_j}-1\big)\Big]\Big|_{x_j=0}
  =\tilde w_{t,j}e^{x_j}\big|_{x_j=0}=\tilde w_{t,j},
\]
i.e.\ the first-order term of $e^{x}-1$ is $x$ itself. This is the algebraic reason why the
linearized loss is the ``delta-weighted return'' of the portfolio.
\end{example}

\begin{example}[Example 1.5: European call, the Greeks]\label{ex:bs}
Let $P$ consist of one standard European call on a non-dividend-paying stock $S$ with maturity $T$
and strike $K$. The Black--Scholes value at time $t$ is
\begin{equation}\label{eq:BS}
  C^{BS}(t,S;r,\sigma)=S\Phi(d_1)-Ke^{-r(T-t)}\Phi(d_2),
\end{equation}
where $\Phi$ is the standard normal df, $r$ the risk-free rate, $\sigma$ the volatility of $S$, and
\[
  d_1=\frac{\log(S/K)+\big(r+\sigma^2/2\big)(T-t)}{\sigma\sqrt{T-t}},
  \qquad d_2=d_1-\sigma\sqrt{T-t}.
\]
In reality interest rates and volatilities fluctuate, so they enter the risk-factor vector:
\[
  \mathbf{Z}_t=(\log S_t,\; r_t,\; \sigma_t),\qquad
  \mathbf{X}_t=\big(\log(S_t/S_{t-1}),\; r_t-r_{t-1},\; \sigma_t-\sigma_{t-1}\big).
\]
Measuring time in years (as in \eqref{eq:BS}) and taking the horizon $\Delta=1/250$ (one day), the
mapping for year $t$ reads
\[
  V_t=f(t,\mathbf{Z}_t)=C^{BS}(t\Delta,S_t;r_t,\sigma_t),
  \qquad\text{so}\qquad
  f_t(t,\mathbf{Z}_t)=C^{BS}_t(t\Delta,S_t;r_t,\sigma_t)\,\Delta .
\]
For a derivative position one usually computes the linearized loss, here with $d=3$ factors:
\[
  L^{\Delta}_{t+1}=-\Big(f_t(t,\mathbf{z}_t)+\sum_{j=1}^{3}f_{z_j}(t,\mathbf{z}_t)X_{t+1,j}\Big)
                 =-\Big(C^{BS}_t\Delta
                        +C^{BS}_S X_{t+1,1}+C^{BS}_r X_{t+1,2}+C^{BS}_\sigma X_{t+1,3}\Big),
\]
in terms of the partial derivatives of the Black--Scholes formula. The Greeks enter:
$C^{BS}_t$ (theta), $C^{BS}_S$ (delta), $C^{BS}_r$ (rho), $C^{BS}_\sigma$ (vega).
\der{$f_t=C^{BS}_t\Delta$}
$V_t=C^{BS}(t\Delta,S_t;r_t,\sigma_t)$; differentiating with respect to the (yearly) calendar
variable $t$ only hits the first argument, $u=t\Delta$, and $\partial u/\partial t=\Delta$. Hence
$f_t=\Delta\cdot\partial_{u}C^{BS}(u,S_t;r_t,\sigma_t)=C^{BS}_t\Delta$.
\selfcheck{the delta hedge of \eqref{eq:BS} is a special case of \eqref{eq:linloss}: for
$K\to0$ (or $S\gg K$) one has $\Phi(d_1)\to1$ so $C^{BS}_S\to1$ and the linearized loss of the option
coincides with that of a linear position in $S$ --- linearisation only ``sees'' the local
behaviour.}
\end{example}

\subsection{The P\&L ansatz}
Summarising, the standard programme is
\[
  L_{t+1}=\ell_{[t]}(\mathbf{X}_{t+1}),\quad \mathbf{X}_{t+1}\in\real^d,\ d\ \text{large};
\]
find a (strongly) stationary model for $\mathbf{X}_{t+1}$; then calculate
$F_{L_{t+1}}(x)=\P(L_{t+1}\le x)$.
\begin{remark}
Finding $F_{L_{t+1}}$ for a general $f$ or $\ell_{[t]}(\cdot)$ is mostly very difficult and rarely
possible, so one aims at \emph{summarising} $F_{L_{t+1}}$. Practice came up with reductions from
$F_{L_{t+1}}$ to so-called \emph{risk measures}.
\end{remark}

\section{Risk Measures}

\subsection{Three approaches to risk measurement}
\begin{itemize}[nosep]
  \item \textbf{Factor sensitivity measures}: give the change in portfolio value for a given
        predetermined change in one underlying risk factor (the Greeks, duration, \dots).
  \item \textbf{Scenario-based risk measures}: consider a number of future scenarios and measure the
        maximum loss of the portfolio under these scenarios (stress tests).
  \item \textbf{Risk measures based on loss distributions}: statistical quantities describing
        characteristics of the loss distribution of the portfolio.
\end{itemize}
They are used for
\begin{itemize}[nosep]
  \item \textbf{determination of risk capital}: the amount of capital needed as a buffer against
        (unexpected) future losses to satisfy a regulator;
  \item \textbf{management tool}: internal limit systems;
  \item \textbf{insurance premia}: a risk measure can be seen as a measure of riskiness of insured
        claims.
\end{itemize}
\emph{Our interpretation}: a risk measure gives the amount of capital that needs to be \emph{added}
to a position with loss $L$ so that the position becomes acceptable to an (internal/external)
regulator.

Risk measures based on loss distributions attempt to quantify the riskiness of a portfolio. The most
popular ones (VaR, ES) describe the right tail of the loss distribution of $L_{t+1}$. Having decided
between conditional/unconditional modelling we simply denote the df of the loss $L=L_{t+1}$ by $F_L$,
$F_L(x)=\P(L\le x)$.

\paragraph{Historical remarks.} 1994: Weatherstone 4 report (J.P.\ Morgan; RiskMetrics).
Oct 1994: founding of ETH RiskLab (what \emph{is} a risk measure?). 1999: Artzner et al.\ (1999):
coherent risk measures.

\subsection{Coherent risk measures}
\paragraph{Notation.} $L^0(\Omega,\mathcal{F},\P)$ is the space of bounded rvs; $M\subseteq
L^0$ is a convex cone ($L_1,L_2\in M\Rightarrow L_1+L_2\in M$; $L\in M,\lambda>0\Rightarrow\lambda L
\in M$). A \textbf{risk measure} is a map
\[
  \rho:M\to\real,\qquad L\mapsto\rho(L).
\]
If the underlying position ($L$) is acceptable, the extra risk capital $\rho(L)$ is needed for
holding $L$. Note that $\rho(L)$ for the (future) loss $L$ is \emph{today's} money, i.e.\ it is
computed from today's portfolio $P$.

\begin{definition}[Axioms of coherence, Artzner et al.\ 1999]
$\rho$ is \textbf{coherent} if
\begin{enumerate}[label=(Ax\arabic*),leftmargin=*]
  \item \textbf{Translation invariance}: $\forall L\in M,\ \forall l\in\real:\
        \rho(L+l)=\rho(L)+l$;
  \item \textbf{Subadditivity}: $\forall L_1,L_2\in M:\
        \rho(L_1+L_2)\le\rho(L_1)+\rho(L_2)$;
  \item \textbf{Positive homogeneity}: $\forall L\in M,\ \lambda>0:\ \rho(\lambda L)=\lambda\rho(L)$;
  \item \textbf{Monotonicity}: $\forall L_1,L_2\in M,\ L_1\le L_2$ (a.s.):\
        $\rho(L_1)\le\rho(L_2)$.
\end{enumerate}
\end{definition}

\paragraph{Questions.}
(Q1) Do coherent risk measures exist? Is there a characterisation?
(Q2) What about existing risk measures in practice --- are they coherent?
(Q3) Does coherence matter?

\paragraph{(Partial) answers.} (A3) It does matter, as the course will show.
(A1) Let $M^1=\{\P:\P$ a probability measure on $(\Omega,\mathcal{F})\}$, let
$\mathcal{P}\subseteq M^1$ be a set of \emph{scenarios}, and
\[
  M_{\mathcal{P}}=\big\{L\in M:\E^{\P}|L|<\infty\ \ \forall\P\in\mathcal{P}\big\}.
\]
One risk measure is
\[
  \rho_{\mathcal{P}}:M_{\mathcal{P}}\to\real,\qquad
  L\mapsto\rho_{\mathcal{P}}(L)=\sup_{\P\in\mathcal{P}}\E^{\P}[L],
\]
i.e.\ \textbf{the worst expected loss over all scenarios}.

\begin{proposition}[Prop.\ 1.7 in MFE (2015); postponed]
\begin{enumerate}[label=\arabic*),leftmargin=*]
  \item For all $\mathcal{P}$: $\rho_{\mathcal{P}}$ is coherent.
  \item Suppose $\Omega=\{\omega_1,\dots,\omega_N\}$ is finite and $M=\{L:\Omega\to\real\}$. Then a
        risk measure $\rho$ is coherent on $M$ if and only if there exists
        $\mathcal{P}\in M^1$ with $\rho=\rho_{\mathcal{P}}$.
\end{enumerate}
\end{proposition}
Roughly speaking, 1) covers \emph{existence} and 2) \emph{uniqueness} of coherent risk measures.

\subsection{Generalised inverses}
\begin{definition}
Let $F$ be a df on $\real$. For all $\alpha\in(0,1)$,
\[
  q_\alpha(F)=F^{\leftarrow}(\alpha)=\inf\{x\in\real:F(x)\ge\alpha\}
\]
denotes the $\alpha$-\emph{quantile} of $F$; $q_{1/2}(F)$ is the median of $F$.
\end{definition}
\begin{remark}
If $F$ is strictly increasing and continuous, then $F^{\leftarrow}=F^{-1}$ (ordinary inverse
function); further properties of generalised inverses are given in the book (MFE 2015,
Section 2.1/Appendix A).
\end{remark}

\subsection{Value-at-Risk and Expected Shortfall}
\begin{definition}[Value-at-Risk]
Given a confidence level $\alpha\in(0,1)$,
\[
  \VaR_\alpha=\VaR_\alpha(L)=F_L^{\leftarrow}(\alpha)
   =\inf\{x\in\real:F_L(x)\ge\alpha\}.
\]
\end{definition}
\begin{definition}[Expected shortfall]
For a loss $L\sim F_L$ with $\E|L|<\infty$, the expected shortfall at level $\alpha\in(0,1)$ is
\[
  \ES_\alpha=\ES_\alpha(L)=\frac{1}{1-\alpha}\int_\alpha^1 \VaR_u(L)\dd u .
\]
\end{definition}

\paragraph{Comments.} VaR is the primary risk measure, ES an alternative. $\VaR_\alpha$ is simply the
$\alpha$-quantile of $F_L$; ES is the average loss when VaR is exceeded and thus conveys information
about the frequency \emph{and} size of large losses --- ES is a ``what-if'' risk measure.
Trivially $\ES_\alpha(L)\ge\VaR_\alpha(L)$ (the average of the upper tail exceeds its lower end), and
ES is also known as Conditional Value-at-Risk (CVaR).

\begin{proposition}[ES is a tail expectation]\label{prop:ES-tails}
If $F_L$ is continuous, then
\[
  \ES_\alpha(L)=\E\big[L\,\big|\,L\ge\VaR_\alpha(L)\big].
\]
If moreover $F_L$ is absolutely continuous with density $f_L$, then
\[
  \ES_\alpha(L)=\frac{1}{1-\alpha}\int_{\VaR_\alpha(L)}^{\infty}x f_L(x)\dd x .
\]
\end{proposition}
\der{Proposition \ref{prop:ES-tails} --- this is the ``(white board)'' computation}
\emph{Step 1 (change of variables).} Substitute $u=F_L(x)$ in the defining integral. Since $F_L$ is
continuous and (for the second part) absolutely continuous on $[\VaR_\alpha(L),\infty)$, we have
$\dd u=f_L(x)\dd x$, and by continuity $F_L(\VaR_\alpha(L))=\alpha$ (no jump can make the df overshoot
$\alpha$ at the quantile). Moreover $u\uparrow1\iff x\uparrow\infty$ and $u\downarrow\alpha\iff
x\downarrow\VaR_\alpha(L)$, so
\[
  \int_\alpha^1 \VaR_u(L)\dd u
   =\int_{\VaR_\alpha(L)}^{\infty} x\,f_L(x)\dd x .
\]
\emph{Step 2 (rewrite as a conditional expectation).} For a continuous $F_L$,
$\P\big(L\ge\VaR_\alpha(L)\big)=1-F_L(\VaR_\alpha(L))=1-\alpha$ (the two-sided version
$\P(L>\VaR_\alpha)=1-\alpha$ holds as well because $\P(L=\VaR_\alpha(L))=0$). Therefore
\[
  \int_{\VaR_\alpha(L)}^{\infty}x f_L(x)\dd x
   =\E\big[L\ind{L\ge\VaR_\alpha(L)}\big]
   =(1-\alpha)\,\E\big[L\,\big|\,L\ge\VaR_\alpha(L)\big],
\]
which gives both statements after dividing by $1-\alpha$. \hfill$\square$
\selfcheck{continuity is essential here. If $F_L$ has an atom at $\VaR_\alpha(L)$ --- say
$F_L(\VaR_\alpha(L)^-)=f<\alpha=F_L(\VaR_\alpha(L))$ --- then the tail event has probability
$1-f>1-\alpha$, and the level-$\alpha$ ES is given by the mixture
\[
  \ES_\alpha(L)=\frac{1}{1-\alpha}\Big(\E\big[L\ind{L\ge\VaR_\alpha(L)}\big]
                +\VaR_\alpha(L)\big(F_L(\VaR_\alpha(L)^-)-\alpha\big)\Big),
\]
which for every integrable $L$ can be rewritten as
\[
  \ES_\alpha(L)=\E\big[L\mid L\ge\VaR_\alpha(L)\big]
   +\frac{\alpha-f}{1-\alpha}\Big(\E\big[L\mid L\ge\VaR_\alpha(L)\big]-\VaR_\alpha(L)\Big)
   \;\ge\;\E\big[L\mid L\ge\VaR_\alpha(L)\big],
\]
with equality if and only if $f=\alpha$. In exercise 1.12 below the df is continuous (normal,
Pareto, $t$), so the simple formula applies.}

\subsection{Practical application}
In banking practice $\VaR=\VaR^M_{\alpha,T}(F_L)$: VaR depends on the significance level $\alpha$, the
holding time $T$, a statistical estimation method $M$, etc. It is a statistical quantity that has to
be estimated from data $\Rightarrow$ statistical uncertainty (see BIS 2012, p.\ 41, Question 8).
Typical choices:
\[
\begin{array}{ll}
  \text{Market risk:} & \alpha=0.99,\ T=10\ \text{d (2 trading weeks)},\\
  \text{Credit and operational risk:} & \alpha=0.999,\ T=1\ \text{a},\\
  \text{Economic capital:} & \alpha=0.9997,\ T=1\ \text{a}.
\end{array}
\]
Overall, estimating such high quantiles requires extreme value theory.
``$\VaR_\alpha(L_1+L_2)$ for $L_i\sim \N(\mu_i,\sigma_i^2)$'' is a meaningless statement unless the
joint df of $(L_1,L_2)$ (marginals \emph{and} dependence) is specified.

\section{Heavy-Tailed Distributions}

\subsection{Definition}
A probability distribution has a \textbf{heavy tail} if its tails decay more slowly than those of an
exponential distribution, i.e.\ the probability of extreme values is significantly higher than for
light-tailed distributions such as the normal. Formally, a distribution with df $F$ and survival
function $\bar F(x)=1-F(x)$ is heavy-tailed if
\[
  \E\big[e^{tX}\big]=\infty\qquad\text{for all }t>0,
\]
i.e.\ the moment generating function (MGF) does not exist for any positive $t$. This implies that the
distribution assigns relatively high probability to large deviations from the mean. Heavy tails are
associated with frequent extreme events/outliers and with slow (e.g.\ power-law) decay of the
survival function, whereas light-tailed distributions (normal, exponential) have rapidly decaying
tails (exponential or faster).
\der{why ``no MGF'' is the same as ``tails heavier than exponential''}
Let $X\ge0$ with survival $\bar F$. Then
\[
  \E\big[e^{tX}\big]=\int_0^\infty e^{tx}f(x)\dd x
    =\underbrace{\big[-e^{tx}\bar F(x)\big]_0^\infty}_{\bar F(0)+\ \text{boundary}}
     +t\int_0^\infty e^{tx}\bar F(x)\dd x
\]
by integration by parts, so $\E[e^{tX}]<\infty$ for some $t>0$ requires $\bar F(x)$ to decay faster
than $e^{-tx}$ for \emph{that} $t$; requiring the failure of finiteness \emph{for every} $t>0$ means
$\bar F\ge e^{-tx}$ infinitely often for every $t>0$ --- exactly ``slower than any exponential''.
For a Pareto tail $\bar F(x)=x^{-1/\xi}$ this holds trivially, and for the Cauchy and the Student-$t$
distribution as well.

\subsection{Example 1: the Pareto distribution}
The Pareto distribution is the classic heavy-tailed model (income, wealth, city sizes, insurance
claims). In the one-parameter (``tail'') parametrisation used in this course,
\[
  \bar F(x)=\P(X>x)=\Big(\frac{x_m}{x}\Big)^{\alpha},\qquad x\ge x_m>0,
\]
with scale $x_m$ (minimum value) and shape/tail index $\alpha>0$. The pdf is
\[
  f(x)=\frac{\alpha x_m^{\alpha}}{x^{\alpha+1}},\qquad x\ge x_m .
\]
The tail decays polynomially (power law), much slower than exponential decay; for small $\alpha$ the
tail is very heavy. \textbf{Moments:} $\E[X^k]$ exists if and only if $k<\alpha$: for $\alpha\le1$ the
mean is infinite, for $\alpha\le2$ the variance is infinite --- a prime example of a distribution with
possibly undefined mean or variance.
\der{moment condition $k<\alpha$}
For $x_m=1$ (the general case is the same up to constants),
\[
  \E[X^k]=\int_1^\infty x^k\,\alpha x^{-\alpha-1}\dd x
         =\alpha\int_1^\infty x^{k-\alpha-1}\dd x
         =\alpha\Big[\frac{x^{k-\alpha}}{k-\alpha}\Big]_1^\infty ,
\]
which is finite iff $k-\alpha<0$, i.e.\ $k<\alpha$ (then $\E[X^k]=\alpha/(\alpha-k)$).

\subsection{Example 2: the Cauchy distribution}
\[
  f(x)=\frac{1}{\pi\gamma}\cdot\frac{1}{1+\big(\frac{x-x_0}{\gamma}\big)^2},\qquad x\in\real,
\]
where $x_0$ is a location and $\gamma>0$ a scale parameter; the standard Cauchy has $x_0=0,\gamma=1$:
$f(x)=1/\big(\pi(1+x^2)\big)$.
Despite its bell shape the Cauchy distribution is heavy-tailed: its survival function decays like
$1/x$, so $\E|X|=\infty$ and $\E[X^2]=\infty$; the mean and variance \textbf{do not exist}. The sample
mean of Cauchy data does not converge to a fixed value as $n\to\infty$ --- it remains erratic. The
Cauchy is the Student-$t$ distribution with 1 degree of freedom.
\der{survival, mean and variance of the standard Cauchy}
The df is $F(x)=\tfrac12+\tfrac1\pi\arctan x$, hence
\[
  \bar F(x)=\tfrac12-\tfrac1\pi\arctan x=\frac1\pi\arctan\frac1x\sim\frac{1}{\pi x}\quad(x\to\infty),
\]
using $\arctan(1/x)=\pi/2-\arctan x$ and $\arctan y\sim y$ as $y\to0$. For the first absolute moment,
\[
  \E|X|=\frac{2}{\pi}\int_0^\infty\frac{x}{1+x^2}\dd x
       =\frac{1}{\pi}\Big[\log(1+x^2)\Big]_0^\infty=\infty
\]
(logarithmic divergence). Moreover
\[
  \E[X^2]=\frac{2}{\pi}\int_0^\infty \frac{x^2}{1+x^2}\dd x
        =\frac{2}{\pi}\int_0^\infty\Big(1-\frac{1}{1+x^2}\Big)\dd x
        =\frac{2}{\pi}\Big[x-\arctan x\Big]_0^\infty=\infty,
\]
and $\E[X]=\int_\real x f(x)\dd x$ is not even well defined as a Lebesgue integral, since
$\E|X|=\infty$. Note that the tail
$\bar F(x)\sim1/(\pi x)$ is a Pareto tail with $\alpha=1$: by the moment condition above the first
moment is exactly at the borderline and diverges.

\subsection{Example 3: the Student-$t$ distribution}
With $\nu$ degrees of freedom,
\[
  f(t)=\frac{\Gamma\big(\frac{\nu+1}{2}\big)}{\sqrt{\nu\pi}\;\Gamma\big(\frac{\nu}{2}\big)}
        \Big(1+\frac{t^2}{\nu}\Big)^{-\frac{\nu+1}{2}},\qquad t\in\real,\ \nu>0 .
\]
For small $\nu$ (e.g.\ $\nu=1,2,3$) the tails are heavy: $f(t)\propto|t|^{-(\nu+1)}$, polynomial
decay, so $1-F_L(x)\sim c_\nu x^{-\nu}$ (l'H\^opital's rule), a power law. Also:
\begin{itemize}[nosep]
  \item if $\nu>1$ the mean exists and equals $0$;
  \item if $\nu>2$ the variance exists and equals $\nu/(\nu-2)$;
  \item for $\nu\le2$ the variance is infinite;
  \item as $\nu\to\infty$ the $t$-distribution approaches the normal distribution (light-tailed).
\end{itemize}
So the $t$-distribution is heavy-tailed for small $\nu$, which makes it useful in robust statistics and
financial modelling.

\der{power-law tail $\bar F(x)\sim c_\nu x^{-\nu}$}
As $|t|\to\infty$,
\[
  f(t)\sim\frac{\Gamma\big(\frac{\nu+1}{2}\big)}{\sqrt{\nu\pi}\,\Gamma\big(\frac\nu2\big)}
        \Big(\frac{t^2}{\nu}\Big)^{-\frac{\nu+1}{2}}
       =\underbrace{\frac{\Gamma\big(\frac{\nu+1}{2}\big)\nu^{\nu/2}}
                       {\sqrt{\nu\pi}\,\Gamma\big(\frac\nu2\big)}}_{=:c}
        \,|t|^{-(\nu+1)} .
\]
Then, by l'H\^opital (both terms $\to0$),
\[
  \lim_{x\to\infty}\frac{\bar F(x)}{x^{-\nu}}
   =\lim_{x\to\infty}\frac{f(x)}{\nu x^{-\nu-1}}
   =\frac{c}{\nu},
\]
so $\bar F(x)\sim c_\nu x^{-\nu}$ with $c_\nu=c/\nu
=\frac{\Gamma(\frac{\nu+1}{2})\nu^{\nu/2-1}}{\sqrt{\pi}\,\Gamma(\frac\nu2)}$.
\der{$\nu/(\nu-2)$ for the variance and $\N(0,1)$ as $\nu\to\infty$}
(i) $X\sim t_\nu$ can be written as $X=Z/\sqrt{V/\nu}$ with $Z\sim\N(0,1)\perp V\sim\chi^2_\nu$;
$\E[X^2]=\E[Z^2]\E[\nu/V]=\nu\,\E[V^{-1}]=\nu/(\nu-2)$ for $\nu>2$
(using $\E[V^{-1}]=1/(\nu-2)$ for $\chi^2_\nu$).
(ii) For the density,
\[
  \Big(1+\frac{t^2}{\nu}\Big)^{-\frac{\nu+1}{2}}
  =\underbrace{\Big[\Big(1+\frac{t^2}{\nu}\Big)^{\nu/2}\Big]^{-1}}_{\to\,e^{-t^2/2}}
   \cdot\underbrace{\Big(1+\frac{t^2}{\nu}\Big)^{-1/2}}_{\to\,1},
\]
and the normalising constant converges by the classical asymptotics
$\Gamma(z+a)/\Gamma(z)\sim z^{a}$: with $z=\nu/2$, $a=1/2$,
\[
  \frac{\Gamma\big(\frac{\nu+1}{2}\big)}{\sqrt{\nu\pi}\,\Gamma\big(\frac\nu2\big)}
  =\frac{\Gamma\big(\frac{\nu}{2}+\frac12\big)}{\Gamma\big(\frac\nu2\big)\sqrt{\nu\pi}}
  \sim\frac{(\nu/2)^{1/2}}{\sqrt{\nu\pi}}=\frac{1}{\sqrt{2\pi}},
\]
so $f_\nu(t)\to\varphi(t)$, the standard normal density.

\subsection{Comparison with light-tailed distributions}
\begin{center}
\begin{tabular}{lll}
\toprule
Distribution & Tail decay & Variance finite?\\
\midrule
Normal $\N(0,1)$ & exponential (fast) & yes\\
Exponential $\Exp(1)$ & exponential & yes\\
Pareto $\Par(\alpha=1.5)$ & power law (slow) & no\\
Cauchy & power law & no\\
$t$-distribution $(\nu=3)$ & power law & yes $(=3)$\\
$t$-distribution $(\nu=1)$ & power law & no\\
\bottomrule
\end{tabular}
\end{center}
Expressed as survival tails (with the density exponent in brackets, as on the slide):
\[
\begin{array}{lll}
\toprule
\text{Distribution} & \bar F(x)\ \big(\text{density }f(x)\big) & \E[X^2]<\infty?\\
\midrule
\N(0,1) & \text{exponential},\ e^{-x^2/2} & \text{yes}\\
\Exp(1) & e^{-x} & \text{yes}\\
\Par(\alpha=1.5) & x^{-1.5}\ \big(f\sim x^{-2.5}\big) & \text{no}\\
\text{Cauchy} & \sim \frac{1}{\pi x}\ \big(f\sim x^{-2}\big) & \text{no}\\
t_\nu,\ \nu=3 & c_3x^{-3}\ \big(f\sim|t|^{-4}\big) & \text{yes }(=3)\\
t_\nu,\ \nu=1 & c_1x^{-1}\ \big(f\sim|t|^{-2}\big) & \text{no}\\
\bottomrule
\end{array}
\]
Only the variance column is compared here; note that the tail decay of the survival function is the
relevant quantity for moments and for EVT, and that it is always one power of $x$ slower than the
density.

\subsection{Applications}
Heavy-tailed distributions are essential wherever extreme events are more common than the normal
model predicts:
\textbf{Finance} (stock returns, insurance claims, market crashes);
\textbf{network traffic} (file sizes, web request delays);
\textbf{environmental data} (flood levels, earthquake magnitudes);
\textbf{actuarial science} (large insurance losses);
\textbf{robust statistics} (to avoid sensitivity to outliers).

\paragraph{Conclusion.} Heavy tails are characterised by slow, typically power-law decay, leading to
a higher probability of extreme values. Unlike light-tailed distributions they may have infinite
moments (mean, variance), which makes classical statistical methods less reliable. Examples: Pareto,
Cauchy, low-degree-of-freedom $t$. Understanding heavy tails is crucial in risk assessment, extreme
value theory, and the modelling of rare but impactful events.

\section{VaR and ES under Heavier Tails: Exercise 1.12}

\begin{example}[1. Normal case]\label{ex:normal}
If $L\sim\N(\mu,\sigma^2)$, then for $\alpha>1/2$
\[
  \VaR_\alpha(L)=\mu+\sigma\Phi^{-1}(\alpha),\qquad
  \ES_\alpha(L)=\mu+\sigma\frac{\varphi\big(\Phi^{-1}(\alpha)\big)}{1-\alpha},
\]
and
\[
  \lim_{\alpha\uparrow1}\frac{\ES_\alpha(L)}{\VaR_\alpha(L)}=1,
\]
where $\varphi(x)=\Phi'(x)=\exp(-x^2/2)/\sqrt{2\pi}$.
\der{VaR, ES and the limit}
Let $Z=(L-\mu)/\sigma\sim\N(0,1)$ and $z_\alpha=\Phi^{-1}(\alpha)>0$ (for $\alpha>1/2$). Since
$F_L$ is continuous and strictly increasing, $F_L^{\leftarrow}=F_L^{-1}$ and
\[
  \VaR_\alpha(L)=\mu+\sigma F_Z^{\leftarrow}(\alpha)=\mu+\sigma z_\alpha .
\]
For ES, use Proposition \ref{prop:ES-tails} and substitute $x=\mu+\sigma z$:
\[
  \ES_\alpha(L)=\frac{1}{1-\alpha}\int_{\mu+\sigma z_\alpha}^{\infty}
        x\,\frac{1}{\sigma}\varphi\Big(\frac{x-\mu}{\sigma}\Big)\dd x
   =\frac{1}{1-\alpha}\int_{z_\alpha}^{\infty}(\mu+\sigma z)\varphi(z)\dd z .
\]
Now $\int_{z_\alpha}^\infty\varphi(z)\dd z=1-\alpha$ (this is the first term $\mu$), and for the second
term use $z\varphi(z)=-\varphi'(z)$:
\[
  \int_{z_\alpha}^{\infty} z\varphi(z)\dd z
   =\Big[-\varphi(z)\Big]_{z_\alpha}^{\infty}=\varphi(z_\alpha),
\]
hence $\ES_\alpha(L)=\mu+\sigma\,\varphi(z_\alpha)/(1-\alpha)$, as claimed.
For the limit, write $z=z_\alpha$ and use Mills' ratio / the standard Gaussian tail asymptotics
$1-\Phi(z)\sim\varphi(z)/z$ as $z\to\infty$:
\[
  \frac{\ES_\alpha(L)}{\VaR_\alpha(L)}
  =\frac{\mu+\sigma\,\varphi(z)/(1-\Phi(z))}{\mu+\sigma z}
  \xrightarrow[z\to\infty]{}\frac{\mu+\sigma z}{\mu+\sigma z}=1
\]
(the numerator satisfies $\varphi(z)/(1-\Phi(z))\to z$; for $\mu<\infty$ the ratio tends to $1$).
So in the normal world ES and VaR are asymptotically indistinguishable in \emph{relative} terms, which
is exactly what fails for heavy tails.
\end{example}

\begin{example}[2. Pareto]\label{ex:pareto}
If $L\sim\Par(1/\xi)$, $0<\xi<1$, i.e.\ $\P(L>x)=x^{-1/\xi}$, $x\ge1$, then
\[
  \VaR_\alpha(L)=(1-\alpha)^{-\xi},\qquad
  \ES_\alpha(L)=\frac{\VaR_\alpha(L)}{1-\xi},\qquad
  \frac{\ES_\alpha(L)}{\VaR_\alpha(L)}=\frac{1}{1-\xi}>1 .
\]
Here the difference between VaR and ES matters.
\der{VaR, ES, ratio}
\emph{VaR.} $F_L(x)=1-x^{-1/\xi}$ for $x\ge1$ is continuous and strictly increasing on $[1,\infty)$, so
solve $1-x^{-1/\xi}=\alpha$: $x^{-1/\xi}=1-\alpha$, i.e.\
$x=(1-\alpha)^{-\xi}=q_\alpha$.
\emph{ES.} Use Proposition \ref{prop:ES-tails} with $f_L(x)=\frac1\xi x^{-1/\xi-1}$ ($x\ge1$):
\[
  \int_{q_\alpha}^{\infty}x f_L(x)\dd x
  =\frac1\xi\int_{q_\alpha}^{\infty}x^{-1/\xi}\dd x
  =\frac1\xi\cdot\frac{q_\alpha^{1-1/\xi}}{1/\xi-1}
  =\frac{q_\alpha^{1-1/\xi}}{1-\xi},
\]
where we used $1-1/\xi<0$ so that the antiderivative vanishes at $+\infty$. Dividing by
$1-\alpha=q_\alpha^{-1/\xi}$:
\[
  \ES_\alpha(L)=\frac{q_\alpha^{1-1/\xi}\cdot q_\alpha^{1/\xi}}{1-\xi}
               =\frac{q_\alpha}{1-\xi}=\frac{\VaR_\alpha(L)}{1-\xi}.
\]
The ratio is $1/(1-\xi)>1$ for every $\alpha$, a constant --- no asymptotics needed.
\end{example}

\begin{example}[3. Student-$t$]\label{ex:t}
If $L\sim t_\nu$ (Student-$t$ on $\nu$ degrees of freedom), i.e.
\[
  f_L(x)=\frac{\Gamma\big(\frac{\nu+1}{2}\big)}{\sqrt{\nu\pi}\Gamma\big(\frac\nu2\big)}
        \Big(1+\frac{x^2}{\nu}\Big)^{-\frac{\nu+1}{2}},\qquad x\in\real,\ \nu>0,
\]
(so that $1-F_L(x)\sim c_\nu x^{-\nu}$, power-law type), then for $\nu>1$
\[
  \lim_{\alpha\uparrow1}\frac{\ES_\alpha(L)}{\VaR_\alpha(L)}=\frac{\nu}{\nu-1}>1 .
\]
In finance often $\nu\in(2,5)$: with $\nu=3$, $\ES_\alpha(L)$ is $50\%$ larger than $\VaR_\alpha(L)$ in
the limit for large $\alpha$.
\der{the limit $\nu/(\nu-1)$ (a general regular-variation computation)}
Suppose (as for the $t_\nu$) that $\bar F_L(x)\sim c x^{-1/\xi}$ with tail index $\xi>0$; then
$1-F_L(\VaR_\alpha(L))=1-\alpha$ gives
\[
  \VaR_\alpha(L)\sim\big(c/(1-\alpha)\big)^{\xi}\quad(\alpha\uparrow1).
\]
Inserting this into the definition,
\begin{align*}
  \ES_\alpha(L)&=\frac{1}{1-\alpha}\int_\alpha^1\VaR_u(L)\dd u
   \sim\frac{1}{1-\alpha}\int_\alpha^1\Big(\frac{c}{1-u}\Big)^{\xi}\dd u\\
   &=\frac{c^{\xi}}{1-\alpha}\Big[\frac{(1-u)^{1-\xi}}{1-\xi}\Big]_{\alpha}^{1}
   =\frac{c^{\xi}(1-\alpha)^{-\xi}}{1-\xi}
   =\frac{\VaR_\alpha(L)}{1-\xi}\quad(\xi<1),
\end{align*}
using $\int_\alpha^1(1-u)^{-\xi}\dd u=\frac{(1-\alpha)^{1-\xi}}{1-\xi}$ for $\xi<1$ (for $\xi\ge1$ the
integral diverges while $\VaR_u$ stays finite, so the two cannot be comparable at all --- in that
regime ES is infinite or the tail index has to be read directly from the exact df). Hence
\[
  \frac{\ES_\alpha(L)}{\VaR_\alpha(L)}\longrightarrow\frac{1}{1-\xi}
  \qquad(\alpha\uparrow1).
\]
For the $t_\nu$ we have $\bar F_L(x)\sim c_\nu x^{-\nu}$, i.e.\ $1/\xi=\nu$, so $\xi=1/\nu$ and the
limit equals $\dfrac{1}{1-1/\nu}=\dfrac{\nu}{\nu-1}$; for \mbox{$\nu=3$} this is $3/2$, i.e.\ $50\%$ more
than VaR (in the limit). By contrast, for the Gaussian tail $\xi\to0$ and the ratio tends to $1$, as in
Example \ref{ex:normal}.
\textbf{Conclusion:} with power-type (heavier) tails it can make a difference which risk measure one
chooses!
\end{example}

\section{Coherence of ES and the Failure of Subadditivity for VaR}

\begin{proposition}[Proposition 6.9 in McNeil et al.\ (2015)]
For all $0<\alpha<1$ and $L\in L^1$ (i.e.\ $\E|L|<\infty$): $\ES_\alpha(L)$ is coherent. In
particular, $\ES_\alpha(L)$ is subadditive.
\end{proposition}
\begin{proof}[Proof (sketch; postponed to future lectures)]
See McNeil et al.\ (2015). The representation $\ES_\alpha=\frac{1}{1-\alpha}\int_\alpha^1\VaR_u\dd u$
shows translation invariance, positive homogeneity and monotonicity immediately; subadditivity is the
hard part and follows from the representation of ES as a supremum of expectations over a suitable set
of scenarios, cf.\ Proposition 1 above.
\end{proof}

VaR always satisfies (Ax1), (Ax3), (Ax4) of coherence, so the only questionable axiom is (Ax2),
subadditivity. \textbf{Statement: in general VaR is not subadditive.}

\subsection{Pareto example: VaR is superadditive}
\begin{example}[Heavy-tailedness, Example 6]
Let $L_1,L_2\stackrel{\text{iid}}{\sim}\Par(1/2)$, i.e.\ $\P(L_1>x)=x^{-1/2}$ for $x\ge1$; note that
$\E[L_1]=\infty$.
\begin{enumerate}[label=\arabic*),leftmargin=*]
  \item One can show (tedious) that
        \[
          \P(L_1+L_2>x)\simeq 2x^{-1/2}-x^{-1},\qquad x\ge2 .
        \]
  \item Deduce that for all $\alpha\in(0,1)$,
        \[
          \VaR_\alpha(L_1+L_2)>\VaR_\alpha(L_1)+\VaR_\alpha(L_2),
        \]
        i.e.\ \textbf{$\VaR_\alpha$ is superadditive}.
\end{enumerate}
\der{the exact tail of the sum, and the superadditivity}
\emph{Step 1 (exact survival of the sum).} With $F_{L_1}(y)=1-y^{-1/2}$ and
$f_{L_1}(y)=\frac12y^{-3/2}$ on $[1,\infty)$, for $x\ge2$,
\[
  \P(L_1+L_2\le x)=\int_1^{x-1}\P\big(L_2\le x-y\big)\,f_{L_1}(y)\dd y
   =\int_1^{x-1}\Big(1-(x-y)^{-1/2}\Big)\frac12y^{-3/2}\dd y .
\]
The first piece integrates directly: $\int_1^{x-1}\frac12y^{-3/2}\dd y=[-y^{-1/2}]_1^{x-1}
=1-(x-1)^{-1/2}$. For the second piece substitute $y=xv$:
\[
  \frac12\int_1^{x-1}(x-y)^{-1/2}y^{-3/2}\dd y
  =\frac{1}{2x}\int_{1/x}^{(x-1)/x}(1-v)^{-1/2}v^{-3/2}\dd v
  =\frac{1}{2x}\Big[-2(1-v)^{1/2}v^{-1/2}\Big]_{1/x}^{(x-1)/x},
\]
because $\frac{\dd}{\dd v}\big[-(1-v)^{1/2}v^{-1/2}\big]\cdot(-2)
=v^{-3/2}(1-v)^{-1/2}$. Evaluating the bracket gives
$\frac{1}{x}\big[(x-1)^{1/2}-(x-1)^{-1/2}\cdot1\big]
=\frac{x-2}{x\sqrt{x-1}}$. Putting the two pieces together,
\[
  \P(L_1+L_2\le x)=1-(x-1)^{-1/2}-\frac{x-2}{x\sqrt{x-1}}
    =1-\frac{2\sqrt{x-1}}{x},
\]
hence the exact tail
\[
  \boxed{\ \P(L_1+L_2>x)=\frac{2\sqrt{x-1}}{x}=2x^{-1/2}\Big(1-\frac1x\Big)^{1/2},\qquad x\ge2.\ }
\]
The form quoted on the slide, $2x^{-1/2}-1/x$, is a leading-order (asymptotic) version of this:
expanding the exact expression, $2x^{-1/2}-x^{-3/2}+O(x^{-5/2})$, the dominant term is
$2x^{-1/2}=2\,\P(L_1>x)$, and only that dominant term is needed below. Note the sanity check at
$x=2$: the exact formula gives $2\sqrt1/2=1$, while $2x^{-1/2}-1/x$ would give $0.914$ --- the sum
satisfies $L_1+L_2\ge2$ with equality only on a null set, so the exact expression is the correct one.
\emph{Step 2 (VaR of the marginals).} $F_{L_1}(x)=1-x^{-1/2}$ for $x\ge1$, so solving
$1-x^{-1/2}=\alpha$ gives $x=(1-\alpha)^{-2}$ and
\[
  \VaR_\alpha(L_1)=\VaR_\alpha(L_2)=(1-\alpha)^{-2},\qquad
  \VaR_\alpha(L_1)+\VaR_\alpha(L_2)=2(1-\alpha)^{-2}.
\]
\emph{Step 3 (VaR of the sum).} Write $p=1-\alpha\downarrow0$ and let $q=\VaR_\alpha(L_1+L_2)$. Then
$q$ solves $\P(L_1+L_2>q)=p$, and for $q$ large the exact tail behaves like $2q^{-1/2}$:
\begin{align*}
  2\sqrt{\frac{q-1}{q^2}}=p
  &\quad\Longrightarrow\quad \frac{4(q-1)}{q^2}=p^2
  \quad\Longrightarrow\quad q^2p^2-4q+4=0\\
  &\quad\Longrightarrow\quad
  q=\frac{2\Big(1+\sqrt{1-p^2}\Big)}{p^2}\sim\frac{4}{p^2}=4(1-\alpha)^{-2},
\end{align*}
where the exact root was obtained by solving the quadratic and discarding the other root,
\[
  \frac{2\Big(1-\sqrt{1-p^2}\Big)}{p^2}<1 .
\]
Comparing with Step 2,
\[
  \VaR_\alpha(L_1+L_2)\sim 4(1-\alpha)^{-2}>2(1-\alpha)^{-2}
   =\VaR_\alpha(L_1)+\VaR_\alpha(L_2),
\]
and indeed the inequality holds exactly for every $\alpha$ (not only asymptotically) because
$2(1+\sqrt{1-p^2})/p^2>2/p^2$. So VaR violates subadditivity already in this simple iid Pareto
example with infinite mean.
\end{example}

\begin{example}[Example 7: exponential variables, all moments finite]\label{ex:exp}
If $L_1,L_2\stackrel{\text{iid}}{\sim}\Exp(1)$ (all moments finite, $\P(L_1>x)=e^{-x}$), then
$\VaR_\alpha$ is superadditive for $\alpha<0.71$.
\der{the threshold $\alpha\approx0.71$}
\emph{Marginals.} $F_{L_1}(x)=1-e^{-x}$ gives
\[
  \VaR_\alpha(L_1)=-\log(1-\alpha)=:t\ \ (\text{so } e^{-t}=1-\alpha),\qquad
  \VaR_\alpha(L_1)+\VaR_\alpha(L_2)=2t .
\]
\emph{Sum.} $L_1+L_2\sim\Gamma(2,1)$ with $F(x)=1-e^{-x}(1+x)$ for $x\ge0$; let $x$ be the
$\alpha$-quantile of the sum, i.e.\ $(1+x)e^{-x}=1-\alpha=e^{-t}$.
\emph{Superadditivity} is equivalent to $x>2t$, which we now translate into a condition on
$\alpha$. Since $x\mapsto(1+x)e^{-x}$ is decreasing in the relevant range, $x>2t$ holds iff
$(1+2t)e^{-2t}>e^{-t}$, i.e.
\[
  1+2t>e^{t}
  \qquad\Longleftrightarrow\qquad t<t^\ast,
\]
where $t^\ast$ solves $e^{t}=1+2t$: numerically $t^\ast\approx1.2564$. Finally
$t=-\log(1-\alpha)<t^\ast\iff \alpha<1-e^{-t^\ast}\approx 1-0.2847=0.7153$. Hence VaR is
superadditive for $\alpha\lesssim0.715$ (and subadditive for larger $\alpha$: in the far tail the
Laplace/exponential behaviour prevails and the relative-order statistics of the sum become
thin-tailed), the number quoted as $0.71$ on the slide. So even with all moments finite,
subadditivity of VaR can fail --- in the body of the distribution.
\end{example}

\subsection{A third counterexample: skewness}
\begin{example}[Exercise 1.16 (Skewness)]
Take $100$ independent bonds ($T=1$ a, nominal value $100$, $2\%$ yearly coupon, $1\%$ default
probability, no recovery). Hence, for $i\in\{1,\dots,100\}$, the loss of bond $i$ is
\[
  L_i=\begin{cases}-2,&\text{with probability }0.99,\\ 100,&\text{with probability }0.01,\end{cases}
\]
where $-2$ is (minus) the coupon received if there is no default and $100$ is the loss of the nominal
if there is a default. (The slide writes the $0.99$ case as ``$2$''; with the loss convention
$L_i=-(V_{i,1}-v_{i,0})$ of \eqref{eq:PL} the sign is negative --- the position earns the coupon.)
Consider the two portfolios
\[
  P_1\ \text{``super diversified''}:\ \sum_{i=1}^{100}L_i,
  \qquad
  P_2\ \text{``highly concentrated''}:\ 100\,L_1 .
\]
Both have the same total nominal exposure, and
\[
  \VaR_{0.95}\Big(\sum_{i=1}^{100}L_i\Big)>\VaR_{0.95}\big(100L_1\big),
\]
i.e.\ VaR rewards the concentrated portfolio and punishes diversification.
\der{the general computation (as hinted at on the ``Further analysis'' slide)}
Write $L_i=-a+(a+b)\tilde L_i$ with $a=2$ (coupon), $b=100$ (nominal) and
$\tilde L_i\sim\text{B}(1,p)$, $p=0.01$: indeed $\tilde L_i=0\Rightarrow L_i=-2$ and
$\tilde L_i=1\Rightarrow L_i=-2+102=100$. \emph{Diversified portfolio.} With
$\tilde L:=\sum_{i=1}^n\tilde L_i\sim\text{B}(n,p)$, $n=100$,
\[
  \sum_{i=1}^n L_i=-na+(a+b)\sum_{i=1}^n\tilde L_i=-na+(a+b)\tilde L .
\]
Using translation invariance (Ax1) and positive homogeneity (Ax3),
\[
  \VaR_\alpha\Big(\sum_{i=1}^n L_i\Big)=-na+(a+b)\,F^{\leftarrow}_{\text{B}(n,p)}(\alpha),
  \qquad
  \VaR_\alpha(nL_1)=-na+(a+b)\,n\,F^{\leftarrow}_{\text{B}(1,p)}(\alpha),
\]
because $\VaR_\alpha(nL_1)=n\VaR_\alpha(L_1)$ and
$\VaR_\alpha(L_1)=\VaR_\alpha(-a+(a+b)\tilde L_1)=-a+(a+b)F^{\leftarrow}_{\text{B}(1,p)}(\alpha)$.
Therefore (the coupon $a$ and the nominal $b$ cancel!)
\[
  \VaR_\alpha\Big(\sum_{i=1}^n L_i\Big)>\VaR_\alpha(nL_1)
  \iff F^{\leftarrow}_{\text{B}(n,p)}(\alpha)>n\,F^{\leftarrow}_{\text{B}(1,p)}(\alpha).
\]
Now $F^{\leftarrow}_{\text{B}(1,p)}(\alpha)=0$ for $\alpha\le1-p$ and $=1$ for $\alpha>1-p$, and
$F^{\leftarrow}_{\text{B}(n,p)}(\alpha)\in\{0,1,\dots,n\}$. Hence
\begin{itemize}[nosep]
  \item for $\alpha\le1-p$: superadditivity holds iff
        $F^{\leftarrow}_{\text{B}(n,p)}(\alpha)>0\iff\alpha>\P(\tilde L=0)=(1-p)^n$;
  \item for $\alpha>1-p$: superadditivity would require
        $F^{\leftarrow}_{\text{B}(n,p)}(\alpha)>n$, which is impossible.
\end{itemize}
So
\[
  \boxed{\ \text{superadditivity}\iff (1-p)^n<\alpha\le1-p\ }
\]
independent of the coupon and the nominal, exactly as claimed on the slide. For $n=100$, $p=0.01$:
$(1-p)^n=0.99^{100}\approx0.366$, so superadditivity holds for $0.366<\alpha\le0.99$. \emph{Numbers.}
Recall $F^{\leftarrow}_{\text{B}(n,p)}(\alpha)$ is the smallest $k$ with $\P(\tilde L\le k)\ge\alpha$.
$(1-\alpha)=(0.99)^{100}\Rightarrow k=0$; $\alpha=0.95\Rightarrow k=3$; $\alpha=0.99\Rightarrow k=4$;
$\alpha=0.995\Rightarrow k=4$ as well, while $F^{\leftarrow}_{\text{B}(1,0.01)}(\alpha)=0$ for
$\alpha\le0.99$ and $=1$ for $\alpha>0.99$. Hence
\[
  \VaR_{0.95}\Big(\sum L_i\Big)=-200+102\cdot3=106,\quad
  \VaR_{0.95}(100L_1)=-200+10200\cdot0=-200,
\]
\[
  \VaR_{0.99}\Big(\sum L_i\Big)=-200+102\cdot4=208,\quad
  \VaR_{0.99}(100L_1)=-200+10200\cdot0=-200,
\]
and at the $99.5\%$ level the concentrated portfolio switches to its catastrophe scenario,
\[
  \VaR_{0.995}(100L_1)=-200+10200=10000,
\]
whereas the diversified one barely moves ($208$). The
concentrated portfolio has a quantile of $-200$ up to the $99\%$ level (it almost surely earns all
coupons) while its $1\%$ tail is catastrophic; the diversified portfolio has a \emph{higher}
quantile ($106$ at $95\%$, $208$ at $99\%$), simply because it almost surely collects some defaults.
This is the ``skewness'' face of the failure of subadditivity.
\end{example}

\subsection{VaR is subadditive in the multivariate normal world}
\begin{example}[Normality, the ``garden of Eden'']
Let $\alpha>1/2$ and
\[
  (L_1,L_2)\sim \N_2(\boldsymbol\mu,\Sigma),\qquad
  \Sigma=\begin{pmatrix}\sigma_1^2&\rho\sigma_1\sigma_2\\ \rho\sigma_1\sigma_2&\sigma_2^2\end{pmatrix},
\]
then $L_1+L_2\sim\N\big(\mu_1+\mu_2,\ \sigma_1^2+\sigma_2^2+2\rho\sigma_1\sigma_2\big)$. Hence, with
$z_\alpha=\Phi^{-1}(\alpha)>0$,
\[
  \VaR_\alpha(L_1+L_2)=\mu_1+\mu_2+\sqrt{\sigma_1^2+\sigma_2^2+2\rho\sigma_1\sigma_2}\;z_\alpha
  \le\mu_1+\mu_2+(\sigma_1+\sigma_2)z_\alpha=\VaR_\alpha(L_1)+\VaR_\alpha(L_2),
\]
so VaR is subadditive in this model.
\der{the Cauchy--Schwarz step}
Cauchy--Schwarz for the inner product $\langle X,Y\rangle=\E[XY]$ gives $|\rho|\le1$; more
elementarily, the quadratic form $\sigma_1^2+\sigma_2^2+2\rho\sigma_1\sigma_2$ is non-negative
definite iff $|\rho|\le1$, and
\[
  \sigma_1^2+\sigma_2^2+2\rho\sigma_1\sigma_2\le\sigma_1^2+\sigma_2^2+2\sigma_1\sigma_2
  =(\sigma_1+\sigma_2)^2
\]
whenever $\rho\le1$. Taking square roots and multiplying by $z_\alpha>0$ gives the claim. Note where
each ingredient is used: $\rho\le1$ (dependence), $z_\alpha>0$ (i.e.\ $\alpha>1/2$), and normality.
\end{example}

Still, one can show that if $L_i\sim\N(\mu_i,\sigma_i^2)$, $i\in\{1,2\}$, then there exists a joint
model for $(L_1,L_2)$ with these marginals under which VaR is \emph{not} subadditive.
Proposition 4.4 of MFE (2015) shows that VaR is coherent for all elliptical models --- the ``garden of
Eden'' of risk management.

\paragraph{Conclusion.} Watch out for superadditivity whenever one of
\begin{itemize}[nosep]
  \item \textbf{heavy-tailedness},
  \item \textbf{skewness}, or
  \item \textbf{(special) dependence}
\end{itemize}
is present (or is statistically hinted at).

\section{Further Reading}

\paragraph{VaR --- badly defined!} ``VaR is the maximum expected loss of a portfolio over a given time
horizon with a certain confidence level'' is mathematically meaningless and potentially misleading:
in no sense is VaR a maximum loss. We can lose more, sometimes much more, depending on the heaviness
of the tail of the loss distribution. The ``bible'' on VaR is the book by Philippe Jorion
(Jorion 2007).

\begin{itemize}[nosep]
  \item On methods for QRM: McNeil, Frey, Embrechts (2015);
  \item on risk management: Crouhy, Galai, Mark (2012);
  \item on VaR: Jorion (2007).
\end{itemize}

\begin{thebibliography}{9}
\bibitem{artzner1999} Artzner, P., Delbaen, F., Eber, J.\,M., Heath, D. (1999),
  \emph{Coherent measures of risk}, Mathematical Finance \textbf{9}, 203--228.
\bibitem{bis2012} BIS (2012), \emph{Fundamental review of the trading book},
  \url{http://www.bis.org/publ/bcbs219.pdf} (2012-02-03).
\bibitem{crouhy2012} Crouhy, M., Galai, D., Mark, R. (2012), \emph{Risk Management}, 2nd ed.,
  New York: McGraw-Hill.
\bibitem{jorion2007} Jorion, P. (2007), \emph{Value at Risk: The New Benchmark for Managing
  Financial Risk}, 3rd ed., New York: McGraw-Hill.
\bibitem{mfe2015} McNeil, A.\,J., Frey, R., Embrechts, P. (2015), \emph{Quantitative Risk
  Management: Concepts, Techniques, Tools}, Princeton University Press.
\end{thebibliography}

\end{document}
